%% ACM TOCE — JARS Quantitative Template
%%
%% A starting point for quantitative papers submitted to ACM Transactions on
%% Computing Education (TOCE), following the APA Journal Article Reporting
%% Standards (JARS).  Replace the placeholder prose with your own.
%%
%% ---------------------------------------------------------------------------
%% REQUIREMENTS
%%
%% * The official ACM `acmart' class, version 2.19 or later.  Overleaf and
%%   current TeX Live installations already provide it.  Do NOT copy acmart.cls
%%   into this directory: a local copy shadows the installed one, and an
%%   out-of-date class silently breaks tagging, citation style, and formatting.
%%
%% * The LuaLaTeX compiler.  PDF tagging comes from the compiler itself, so
%%   pdfLaTeX will not produce an accessible PDF.  Overleaf's default LuaLaTeX
%%   is sufficient.
%%
%% TOCE requires submissions to be screen readable; untagged PDFs may be
%% returned before review.  See:
%% https://github.com/acm-toce/documentation/wiki/How-do-I-make-my-PDF-accessible
%% ---------------------------------------------------------------------------

%% \DocumentMetadata must be the very first command in the file.
\DocumentMetadata{tagging=on,
   tagging-setup={math/setup=mathml-SE},
   pdfstandard={A-4f,UA-2},
   lang=en-US
}

%% `anonymous' and `review' are what you want while under review: TOCE uses
%% double-anonymous review.  For the accepted version, drop both options.
\documentclass[acmsmall,screen,anonymous,review]{acmart}

\usepackage{multirow}
\usepackage{listings}
\usepackage{subcaption}

\AtBeginDocument{%
  \providecommand\BibTeX{{Bib\TeX}}}

%% Treat the first row of every table as a header row, so screen readers can
%% associate data cells with their headings.  Override per table as needed.
\tagpdfsetup{table/header-rows={1}}

%% Rights management.  These are SAMPLE values; ACM sends you the real ones
%% when you complete the rights form after acceptance.
\setcopyright{acmlicensed}
\copyrightyear{2026}
\acmYear{2026}
\acmDOI{XXXXXXX.XXXXXXX}

%% Journal information.  Leave these placeholders alone for submission.
\acmJournal{TOCE}
\acmVolume{0}
\acmNumber{0}
\acmArticle{0}
\acmMonth{0}

%% TOCE requires [Author Year] citations (journal policy, May 2026).
\citestyle{acmauthoryear}

\begin{document}


%From the APA JARS Specification
%Identify main variables and theoretical issues under investigation and the relationships between them.
%Identify the populations studied.
\title[JARS Quantitative Template for ACM TOCE]{The JARS Quantitative Template for ACM TOCE}

%% Author information.  The `anonymous' option above hides all of this in the
%% compiled PDF, which is what you want under review.  Fill it in anyway: it
%% appears once you remove `anonymous' for the accepted version.
%%
%% Anonymizing the author block is NOT sufficient on its own.  Work through the
%% anonymization checklist in the TOCE Author Guidelines:
%% https://dl.acm.org/journal/toce/author-guidelines
\author{First Author}
\email{first.author@example.edu}
\affiliation{%
  \institution{First Author's Institution}
  \city{City}
  \country{Country}}

\author{Second Author}
\email{second.author@example.edu}
\affiliation{%
  \institution{Second Author's Institution}
  \city{City}
  \country{Country}}

\renewcommand{\shortauthors}{Author et al.}

%% Structured-abstract headings.  Approach adapted from:
%% https://tex.stackexchange.com/questions/74371/subdividing-structured-abstracts
\newcommand{\AbstractCategory}[1]{%
  \par\addvspace{.5\baselineskip}% adjust to suit
  \noindent\textbf{#1}\quad\ignorespaces
}

%
% The abstract is a short summary of the work to be presented in the article.
\begin{abstract}
\AbstractCategory{Objectives}State the problem under investigation, including main hypotheses.


\AbstractCategory{Participants}
Describe participants, specifying their pertinent characteristics for the study. Participants are described in greater detail in the body of the paper.

\AbstractCategory{Study Methods} Describe the study method, including the research design (e.g., experiment, observational study), sample size, materials used (e.g., instruments, apparatus), outcome measures, and data-gathering procedures, including a brief description of the source of any secondary data. If the study is a secondary data analysis, so indicate.

\AbstractCategory{Findings} Report findings, including effect sizes and confidence intervals or statistical significance levels.

\AbstractCategory{Conclusions} State conclusions, beyond just results, and report the implications or applications.

\end{abstract}



%
% The code below is generated by the tool at http://dl.acm.org/ccs.cfm.
% Please copy and paste the code instead of the example below.


\begin{CCSXML}
<ccs2012>
<concept>
<concept_id>10002950.10003705.10003708</concept_id>
<concept_desc>Mathematics of computing~Statistical software</concept_desc>
<concept_significance>500</concept_significance>
</concept>
<concept>
<concept_id>10010405.10010489</concept_id>
<concept_desc>Applied computing~Education</concept_desc>
<concept_significance>500</concept_significance>
</concept>
<concept>
<concept_id>10003120.10003121.10011748</concept_id>
<concept_desc>Human-centered computing~Empirical studies in HCI</concept_desc>
<concept_significance>300</concept_significance>
</concept>
<concept>
<concept_id>10011007.10011006.10011050.10011017</concept_id>
<concept_desc>Software and its engineering~Domain specific languages</concept_desc>
<concept_significance>300</concept_significance>
</concept>
</ccs2012>
\end{CCSXML}

\ccsdesc[500]{Mathematics of computing~Statistical software}
\ccsdesc[500]{Applied computing~Education}
\ccsdesc[300]{Human-centered computing~Empirical studies in HCI}
\ccsdesc[300]{Software and its engineering~Domain specific languages}


%
% Keywords. The author(s) should pick words that accurately describe the work being
% presented. Separate the keywords with commas.
\keywords{comma, separated}


%
% This command processes the author and affiliation and title information and builds
% the first part of the formatted document.
\maketitle

%%Optional
\textbf{Author's Note:} Provide acknowledgment and explanation of any special circumstances, including:

\begin{itemize}
    \item Registration information if the study has been registered
    \item Use of data also appearing in previous publications
    \item Prior reporting of the fundamental data in dissertations or conference papers
    \item Sources of funding or other support
    \item Relationships or affiliations that may be perceived as conflicts of interest
    \item Previous (or current) affiliation of authors if different from location where the study  was conducted
    \item Your positionality with respect to this study, including how your background and orientation may have shaped your choice of theory, study design, sampling frame, measures, and interpretation (JARS--REC). Anonymize this statement while the paper is under review.
    \item Contact information for the corresponding author
    \item Additional information of importance to the reader that may not be appropriately included in other sections of the paper
    \item TOCE uses a double-anonymous review process. Omit author information and all identifying information in papers submitted for review. Please see \href{https://dl.acm.org/journal/toce/author-guidelines}{TOCE’s Author Guidelines page} for advice on anonymizing your paper.
\end{itemize}


\section{Introduction}

ACM TOCE strongly encourages authors to submit papers that adhere to the \href{https://apastyle.apa.org/jars/}{American Psychological Association’s Journal Article Reporting Standards}  (JARS). Some papers that lie within the scope of ACM TOCE are not a good fit for these standards; we welcome such submissions. Whether or not your paper is a good fit for the standards, they can provide useful guidance to help make your paper more accessible to reviewers and the TOCE audience. Here, we provide a JARS template for \href{https://apastyle.apa.org/jars/quant-table-1.pdf}{general quantitative research designs}. Use \href{https://apastyle.apa.org/jars/jars-quant-participant-flowchart.pdf}{this flowchart} to identify the JARS reporting guidelines for your specific quantitative design, including \href{https://apastyle.apa.org/jars/quant-table-2.pdf}{Experimental Designs}, \href{https://apastyle.apa.org/jars/quant-table-2a.pdf}{Random Assignment}, \href{https://apastyle.apa.org/jars/quant-table-2b.pdf}{Non-Random Assignment}, \href{https://apastyle.apa.org/jars/quant-table-2c.pdf}{Clinical Trials}, \href{https://apastyle.apa.org/jars/quant-table-3.pdf}{Non-Experimental Designs}, Special Designs (\href{https://apastyle.apa.org/jars/quant-table-4.pdf}{Longitudinal Studies}, \href{https://apastyle.apa.org/jars/quant-table-5.pdf}{N-of-1 Studies}, \href{https://apastyle.apa.org/jars/quant-table-6.pdf}{Replication Studies}), Analytic Methods (\href{https://apastyle.apa.org/jars/quant-table-7.pdf}{Structural Equation Modeling}, \href{https://apastyle.apa.org/jars/quant-table-8.pdf}{Bayesian Statistics}), and \href{https://apastyle.apa.org/jars/quant-table-9.pdf}{Meta-Analyses}. The reporting standards for these more specific quantitative designs can be readily integrated into this template. See also the TOCE JARS templates for \href{https://apastyle.apa.org/jars/qualitative}{qualitative} and \href{https://apastyle.apa.org/jars/mixed-methods}{mixed} research designs.


Note that ACM provides its own article templates for \href{https://www.acm.org/binaries/content/assets/publications/taps/acm_submission_template.docx}{Microsoft Word} and \href{https://portalparts.acm.org/hippo/latex_templates/acmart-primary.zip}{LaTeX}. Download whichever matches your workflow and consult it as you write: it covers formatting details this template does not. If you use LaTeX, note that TOCE requires \texttt{acmart} version 2.19 or later, which this template assumes.

Citations in TOCE use the \verb|[Author Year]| style, which this template enables with \verb|\citestyle{acmauthoryear}| in the preamble. A citation renders like this~\cite{Bak2017}; add your own entries to \texttt{bibliography.bib}. Before you submit, verify that every reference you cite actually exists and is described correctly --- TOCE requires authors to check their bibliographies against the digital library record for each citation.

Whatever design your study uses, the APA \href{https://apastyle.apa.org/jars/race-ethnicity-culture}{Journal Article Reporting Standards for Race, Ethnicity, and Culture} (JARS--REC, added in 2023) apply as well. JARS--REC is not a fourth design type; it is a cross-cutting standard describing how to report race, ethnicity, and culture in any study, and \href{https://apastyle.apa.org/jars/rec-table-1.pdf}{REC Table 1} sets out what to include section by section. Prompts drawn from it appear at the relevant points below.
\subsection{Problem}

State the importance of the problem, including theoretical or practical implications.

\subsection{Review of Relevant Scholarship}

Provide a succinct review of relevant scholarship, including: 

\begin{itemize}
    \item Relation to previous work
    \item Differences between the current report and earlier reports if some aspects of this study have been reported on previously
\end{itemize}

\subsection{Hypothesis, Aims, and Objectives}

State specific hypotheses, aims, and objectives, including:
\begin{itemize}
    \item theories or other means used to derive hypotheses
    \item primary and secondary hypotheses
    \item other planned analyses
\end{itemize}

State how hypotheses and research design relate to one another.

\section{Method}

\subsection{Inclusion and Exclusion}
Report inclusion and exclusion criteria, including any restrictions based on demographic characteristics.

\subsection{Participant Characteristics}
Report major demographic characteristics (e.g., age, sex, ethnicity, socioeconomic status) and important topic-specific characteristics (e.g., achievement level in studies of educational interventions).

%% JARS-Quant Table 1 also asks for nonhuman-animal reporting details here.
%% TOCE leaves this commented out: it is very rarely relevant to computing
%% education research.  Uncomment it if your study does involve animal subjects.
%In the case of animal research, report the genus, species, and strain number or other specific identification, such as the name and location of the supplier and the stock designation. Give the number of animals and the animals’ sex, age, weight, physiological condition, genetic modification status, genotype, health–immune status, drug or test naïveté, and previous procedures to which the animal may have been subjected.


When you report race, ethnicity, or culture, JARS--REC asks for more than a category count:
\begin{itemize}
    \item Describe the racial and ethnic composition of the sample in as much detail as space allows, including participants' country of origin and language use. If you offered an open-ended ``self-described'' option, report the responses you received.
    \item Give the wording of the questions you used to elicit race, ethnicity, culture, country of origin, and languages spoken, along with the response options offered. If law or policy in the country of data collection prevented you from asking, say so explicitly.
    \item Explain how multiracial and multiethnic responses were handled, and justify any aggregation across identities rather than treating aggregation as the default.
    \item Use racial and ethnic terms that match local usage and that participants themselves would recognize, and say why you chose them. Terminology shifts over time and varies across communities.
\end{itemize}

\subsection{Sampling Procedures}

\begin{itemize}
    \item Describe procedures for selecting participants, including:
    \begin{itemize}
        \item Sampling method if a systematic sampling plan was implemented
        \item Percentage of sample approached that actually participated
        \item Whether self-selection into the study occurred (either by individuals or by units, such as schools or clinics)
    \end{itemize}
    \item Describe settings and locations where data were collected as well as dates of data collection.
    \item Describe agreements and payments made to participants.
    \item Describe institutional review board agreements, ethical standards met, and safety monitoring.
\end{itemize}

\subsection{Sample Size, Power, and Precision}
Describe the sample size, power, and precision, including:
\begin{itemize}
    \item Intended sample size
    \item Achieved sample size, if different from the intended sample size
    \item Determination of sample size, including power analysis, or methods used to determine precision of parameter estimates and explanation of any interim analyses and stopping rules employed
\end{itemize}

\subsection{Measures and Covariates}
Define all primary and secondary measures and covariates, including measures collected but not included in the report.

\subsection{Data Collection}
Describe methods used to collect data.

\subsection{Quality of Measurements}
Describe methods used to enhance the quality of measurements, including:
\begin{itemize}
    \item training and reliability of data collectors
    \item use of multiple observations
\end{itemize}

\subsection{Instrumentation}
Provide information on validated or ad hoc instruments created for individual studies, for individual studies (e.g., psychometric and biometric properties)

\subsection{Masking}
Report whether participants, those administering the experimental manipulations, and those assessing the outcomes were aware of condition assignments. If masking took place, provide a statement regarding how it was accomplished and whether and how the success of masking was evaluated. 

\subsection{Psychometrics}
\begin{itemize}
    \item Estimate and report values of reliability coefficients for the scores analyzed (i.e., the researcher’s sample), if possible. 
    \item Provide estimates of convergent and discriminant validity where relevant.
    \item Report estimates related to the reliability of measures, including: 
    \begin{itemize}
        \item interrater reliability for subjectively scored measures and ratings
        \item test--retest coefficients in longitudinal studies in which the retest interval corresponds to the measurement schedule used in the study
        \item internal consistency coefficients for composite scales in which these indices are appropriate for understanding the nature of the instruments being used in the study
    \end{itemize}
    \item Report the basic demographic characteristics of other samples if reporting reliability or validity coefficients from those samples, such as those described in test manuals or in norming information for the instrument.
\end{itemize}

\subsection{Conditions and Design}

\begin{itemize}
    \item State whether conditions were manipulated or naturally observed. Report the type of design as per the JARS–Quant tables:
    \begin{itemize}
        \item experimental manipulation with participants randomized: \href{https://apastyle.apa.org/jars/quant-table-2.pdf}{Table 2} and \href{https://apastyle.apa.org/jars/quant-table-2a.pdf}{Module A}
        \item experimental manipulation without randomization: \href{https://apastyle.apa.org/jars/quant-table-2.pdf}{Table 2} and \href{https://apastyle.apa.org/jars/quant-table-2b.pdf}{Module B}
        \item clinical trial with randomization: \href{https://apastyle.apa.org/jars/quant-table-2.pdf}{Table 2} and Modules \href{https://apastyle.apa.org/jars/quant-table-2a.pdf}{A} and \href{https://apastyle.apa.org/jars/quant-table-2c.pdf}{C}
        \item clinical trial without randomization: \href{https://apastyle.apa.org/jars/quant-table-2.pdf}{Table 2} and Modules \href{https://apastyle.apa.org/jars/quant-table-2b.pdf}{B} and \href{https://apastyle.apa.org/jars/quant-table-2c.pdf}{C}
        \item nonexperimental design (i.e., no experimental manipulation): observational design, epidemiological design, natural history, and so forth (single-group designs or multiple-group comparisons): \href{https://apastyle.apa.org/jars/quant-table-3.pdf}{Table 3}
        \item longitudinal design: \href{https://apastyle.apa.org/jars/quant-table-4.pdf}{Table 4}
        \item N-of-1 studies: \href{https://apastyle.apa.org/jars/quant-table-5.pdf}{Table 5}
        \item replications: \href{https://apastyle.apa.org/jars/quant-table-6.pdf}{Table 6}
    \end{itemize}

    \item Report the common name given to designs not currently covered in JARS–Quant.
\end{itemize}

\subsection{Data Diagnostics}

Describe planned data diagnostics, including:
\begin{itemize}
    \item criteria for post-data-collection exclusion of participants, if any
    \item criteria for deciding when to infer missing data, and methods used for imputation of missing data
    \item definition and processing of statistical outliers
    \item analyses of data distributions
    \item data transformations to be used, if any
\end{itemize}

\subsection{Analytic Strategy}

Describe the analytic strategy for inferential statistics and protection against experiment-wise error for:
\begin{itemize}
    \item primary hypotheses
    \item secondary hypotheses
    \item exploratory hypotheses
\end{itemize}

\section{Results}

\subsection{Participant Flow}
Report the flow of participants, including:
\begin{itemize}
    \item total number of participants in each group at each stage of the study
    \item flow of participants through each stage of the study (include figure depicting flow, when possible; see the \href{https://apastyle.apa.org/jars/jars-quant-participant-flowchart.pdf}{JARS–Quant Participant Flowchart})
\end{itemize}

\subsection{Recruitment}
Provide dates defining the periods of recruitment and repeated measures or follow-up.

\subsection{Statistics and Data Analysis}
Provide information detailing the statistical and data-analytic methods used, including:

\begin{itemize}
    \item Missing data, including:
    \begin{itemize}
        \item frequency or percentages of missing data
        \item empirical evidence and/or theoretical arguments for the causes of data that are missing---for example, missing completely at random (MCAR), missing at random (MAR), or missing not at random (MNAR)
        \item methods actually used for addressing missing data, if any
    \end{itemize}
    \item Descriptions of each primary and secondary outcome, including the total sample and each subgroup, that includes the number of cases, cell means, standard deviations, and other measures that characterize the data used
    \item Inferential statistics, including:
    \begin{itemize}
        \item results of all inferential tests conducted, including exact \emph{p} values if null hypothesis significance testing (NHST) methods were used, and reporting the minimally sufficient set of statistics (e.g., \emph{df}s, mean square [MS] effect, MS error) needed to construct the tests
        \item effect-size estimates and confidence intervals on estimates that correspond to each inferential test conducted, when possible
        \item clear differentiation between primary hypotheses and their tests--estimates, secondary hypotheses and their tests--estimates, and exploratory hypotheses and their tests--estimates
    \end{itemize}
    \item Complex data analyses---for example, structural equation modeling analyses (see also \href{https://apastyle.apa.org/jars/quant-table-7.pdf}{Table 7}), hierarchical linear models, factor analysis, multivariate analyses, and so forth, including:
    \begin{itemize}
        \item details of the models estimated
        \item associated variance--covariance (or correlation) matrix or matrices
        \item identification of the statistical software used to run the analyses (e.g., SAS PROC GLM or the particular R package)
    \end{itemize}
    \item Estimation problems (e.g., failure to converge, bad solution spaces), regression diagnostics, or analytic anomalies that were detected, and the solutions to those problems
    \item Other data analyses performed, including adjusted analyses, if performed, indicating those that were planned and those that were not planned (though not necessarily in the level of detail of primary analyses)
\end{itemize}

Report any problems with statistical assumptions and/or data distributions that could affect the validity of findings.


\section{Discussion}
\subsection{Support of Original Hypotheses}
Provide a statement of support or nonsupport for all hypotheses, whether primary or secondary, including:
\begin{itemize}
    \item distinction by primary and secondary hypotheses
    \item discussion of the implications of exploratory analyses in terms of both substantive findings and error rates that may be uncontrolled
\end{itemize}

\subsection{Similarity of Results}
Discuss similarities and differences between reported results and the work of others.

\subsection{Interpretation}

Provide an interpretation of the results, taking into account:

\begin{itemize}
    \item Sources of potential bias and threats to internal and statistical validity
    \item Imprecision of measurement protocols
    \item Overall number of tests or overlap among tests
    \item Adequacy of sample sizes and sampling validity
\end{itemize}


\subsection{Generalizability}

Discuss generalizability (external validity) of the findings, taking into account:

JARS--REC additionally asks you to state the constraints on generality: name the populations to which the findings are and are not expected to apply, and avoid implying that a sample drawn from one racial, ethnic, or cultural group generalizes by default.

\begin{itemize}
    \item target population (sampling validity)
    \item other contextual issues (setting, measurement, time; ecological validity)
\end{itemize}

\subsection{Implications}

Discuss implications for future research, program, or policy.

%% acmart's `acks' environment sets the heading for you AND suppresses the whole
%% block while the `anonymous' option is in force -- which a plain \section would
%% not do.  Use it, and your acknowledgments cannot leak into the review copy.
\begin{acks}
Acknowledgments are placed before the references. Add information about grants, awards, or other types of funding you received to support this research, and thank those who contributed to the work without meeting the criteria for authorship.
\end{acks}

\bibliographystyle{ACM-Reference-Format}
\bibliography{bibliography.bib}

\appendix


\end{document}
